\documentclass[11pt]{article}
\usepackage[margin=0.82in]{geometry}
\usepackage{booktabs}
\usepackage{hyperref}

\title{HAOS-IIP Magnon Address Stability: Testable Signatures}
\author{Tomislav Rupi\'c \\ Independent Researcher}
\date{May 2026}

\begin{document}
\maketitle

\section*{Scope}

This note creates the parallel spin-sector prediction layer. It is not a new
spin-wave theory, not an external magnon proof of HAOS-IIP, and not a claim
that ferrons and magnons share a literal substrate.

\section*{Calibration and Boundary}

The committed calibration source is Materials Bridge Line A: ferron persistence
proxy \texttt{0.957192}, mean spectral recoverability \texttt{0.928455},
\texttt{15/15} target-window peaks near \texttt{3.13 THz}, and raw STFT-grid
status \texttt{NO\_STFT\_DATA\_FOUND}.

A user-provided alpha-Fe$_2$O$_3$ magnon audit anchor reports persistence proxy
\texttt{0.902} and raw public STFT-grid boundary
\texttt{NO\_PUBLIC\_STFT\_GRID\_FOUND}. The corresponding audit artifact was not
found in this repo at note creation time, so \texttt{0.902} is recorded as a
user-provided anchor until the audit files are committed.

\section*{Candidate External Anchors}

Candidate audit targets include Li et al., ``Anisotropic Coherent Propagation of
Sub-Terahertz Magnons in Altermagnetic alpha-Fe$_2$O$_3$,'' Advanced Optical
Materials 2026; Yang et al., ``Spin-orbit torque manipulation of
sub-terahertz magnons in antiferromagnetic alpha-Fe$_2$O$_3$,'' Nature
Communications 2024; Sun et al., ``Observation of Chiral Magnon Band Splitting
in Altermagnetic Hematite,'' Physical Review Letters 2025; Choe et al.,
``Long-lived zone-boundary magnons in an antiferromagnet,'' Nature
Communications 2025; and Yuan et al., ``Inelastic Neutron Scattering Study of
Field Dependence of Magnetic Excitations in CoTiO$_3$,'' Physical Review B
2024. These references motivate target datasets; they do not validate HAOS
metrics by themselves.

\section*{Field-Facing Predictions}

\begin{center}
\begin{tabular}{p{0.24\linewidth}p{0.35\linewidth}p{0.31\linewidth}}
\toprule
Prediction & Testable readout & Falsifier \\
\midrule
Crystal-orientation-dependent narrowband magnon address &
Sub-THz exchange magnons should show highest coherence and narrowest linewidth
only for specific wavevector / crystal-axis alignments. &
If orientation scans show isotropic broadband response, or aligned records do
not outperform misaligned controls, the prediction fails. \\
\addlinespace
Thickness / boundary-dependent group velocity &
Magnon packet velocity should scale with thickness, confinement, or boundary
condition while target-window recoverability remains above \texttt{0.70} before
visible propagation failure. &
If controlled thickness or boundary series show no systematic velocity scaling
and no distance-ordered delay beyond uncertainty, the prediction fails. \\
\addlinespace
Target-address recovery after perturbation &
After a secondary pulse, field quench, strain perturbation, or spin-orbit torque
below damage threshold, the readout should return to the original target
frequency window with recoverability at least \texttt{0.85}. &
If perturbation repeatedly causes complete thermalization or a different
dominant mode with no return to the original address, the prediction fails. \\
\addlinespace
Hybrid magnon-phonon / cross-sector locking &
In vdW magnetic hybrids or coupled heterostructures, compatible magnon and
phonon / polar modes should produce enhanced lifetime, cleaner linewidth, or
more directional propagation than uncoupled controls. &
If hybridization gives no lifetime, linewidth, directionality, or recoverability
benefit beyond uncoupled controls, the prediction fails. \\
\bottomrule
\end{tabular}
\end{center}

\section*{Boundary}

These are field-facing diagnostic candidates. They do not claim new physics, a
universal harmonic grid, or a replacement for spin-wave theory. The next
required step is to commit the alpha-Fe$_2$O$_3$ audit artifact or build a new
magnon data bridge that loads public spectra / traces and runs the same bounded
recoverability audit used for ferrons.

\end{document}
